% Responsável: TODO(grupo)
\chapter{Introdução}
\label{cap:introducao}

% Contexto do problema: falhas características em máquinas rotativas e a
% distinção entre DETECTAR que algo mudou e DIAGNOSTICAR o que mudou.
% Fechar com a estrutura do relatório.

Rolamentos estão entre os componentes que mais falham em máquinas rotativas, e
uma falha não detectada costuma terminar em parada não programada. O
monitoramento por vibração permite acompanhar a condição da máquina em operação:
um defeito localizado gera impactos periódicos, cuja taxa depende da geometria do
rolamento e da rotação, e que deixam no sinal uma assinatura reconhecível
~\cite{randall2011livro,randall2011tutorial}.

Duas perguntas diferentes podem ser feitas a esse sinal. A primeira é de
\emph{detecção}: o comportamento da máquina mudou em relação ao normal? A
segunda é de \emph{diagnóstico}: se mudou, qual é o defeito? As duas exigem
dados diferentes. Para detectar, basta conhecer o comportamento normal, que é o
único disponível numa planta antes da primeira falha. Para diagnosticar, é
preciso ter exemplos de cada tipo de defeito.

Este trabalho trata as duas perguntas com os dados de rolamentos da
\emph{Case Western Reserve University} (CWRU)~\cite{cwru}, a referência mais usada
na área, mas que exige cuidado: as gravações variam muito de qualidade, e parte
delas não mostra a assinatura do defeito declarado~\cite{smith2015}. Por isso a
ênfase está no protocolo de avaliação --- divisão por condição de carga, nenhum
vazamento entre treino e teste, detector treinado só com dados normais --- e na
verificação física dos resultados, comparando o que os modelos indicam com as
frequências de defeito previstas pela teoria.

O \autoref{cap:dataset} descreve o conjunto de dados e a auditoria das gravações.
O \autoref{cap:fundamentacao} apresenta a física dos defeitos e as técnicas
utilizadas, e o \autoref{cap:metodologia}, o pré-processamento, o protocolo e as
features. Os resultados estão no \autoref{cap:resultados}, discutidos
criticamente no \autoref{cap:discussao}, e o \autoref{cap:conclusao} conclui.

\section{Objetivo}
\label{sec:objetivo}

O objetivo é avaliar, em dados reais de vibração do CWRU, duas tarefas tratadas
separadamente:
\begin{enumerate}
  \item \textbf{Detecção de anomalia} sem exemplos de falha: treinar detectores
  \emph{somente} com sinais de rolamento normal e medir a taxa de verdadeiros e de
  falsos positivos ao aplicá-los a uma condição de carga não vista;
  \item \textbf{Diagnóstico do tipo de falha}: classificar as janelas de sinal em
  normal, defeito na pista interna, na pista externa ou na esfera, também numa
  carga não vista no treino.
\end{enumerate}
Em ambas, o interesse está tanto no desempenho quanto em onde e por que os
modelos erram, e em se as decisões se apoiam na física do defeito.
